\documentclass[10pt,twocoloumn]{article}

\usepackage{amsmath}
\usepackage{amsthm}
\usepackage{kida}

\usepackage{float}
\usepackage{svg}
\usepackage{graphicx}
\newcommand{\addsvg}[4][1.0]{
    \begin{figure}[H]
        \centering
        \includesvg[width=#1\textwidth]{#2}
        \caption{#4}
        \label{fig:#3}
    \end{figure}
}

\title{Gravitation 1 Assignment 1}
\author{Amir Hossein Ebrahimnezhad - 404461001}
\date{due: \today}

\begin{document}
\maketitle
\section*{Problem 1}
\subsection*{Problem Statement}
For Newtonian gravity the potential is related to the mass density by Poisson's equation
\begin{equation}
\nabla^2\Phi(\vec{x},t) = 4\pi G \rho(\vec x , t)
\end{equation}
Before Special relativity, Galilean transformation was used to relate the space and time coordinates of an event as observed in two different inertial reference.

\begin{itemize}
    \item Show that the Poisson's equation is invariant under a Galilean transformation.
    \item Then show that Poisson's equation is not invariant under a Lorentz transformation.
\end{itemize}
\subsection*{Solution}
Let's consider to coordinate systems $S$ and $S'$ with relative velocity $v$. Consider the Poisson's equation in $S$ coordinate as:

\begin{equation}
\left(\frac{\partial^2}{\partial x^2} + \frac{\partial^2}{\partial y^2} + \frac{\partial^2}{\partial z^2}\right)\Phi  = \alpha \rho
\end{equation}

where $\alpha = 4\pi G$. Galilean transformation imposes the following rules of transformation from $S$ to $S'$:

\begin{align}
x'&= x +  vt\\
y'&=y\\
z'&=z \\
t'&=t
\end{align}

we therefore expand the Poisson's equation and write in terms of the $S'$ coordinates:

\begin{equation}
\partial_{x'}\left(\partial_{x'}\Phi \partial_xx' \right)\partial_xx' + \partial_{y'}\left(\partial_{y'}\Phi \partial_yy' \right)\partial_yy' +\partial_{z'}\left(\partial_{z'}\Phi \partial_zz' \right)\partial_zz'
\end{equation}

now we know by the rules that:
\begin{equation}
    \partial_i i' = 1
\end{equation}
where $i=x,y,z$ and $i'=x',y',z'$

Therefore:
\begin{equation}
\nabla^2 = \nabla'^2
\end{equation}
Now the Poisson's equation becomes:
\begin{equation}
\nabla'^2 \Phi(\vec x',t') = 4\pi G\rho(\vec x' , t'), \ \qed
\end{equation}

Following this, we know that the transformation rules imposed by Lorentz transformations are:
\begin{align}
t' &= \gamma(t - v x)\\
x' &= \gamma(x - v t)\\
y' &= y\\
z' &= z
\end{align}
where $\gamma \equiv (1-v^2)^{-1/2}$ and we work in units with $c=1$.

Using the chain rule we express derivatives in the unprimed frame in terms of primed derivatives:
\begin{align}
\partial_x &= \frac{\partial x'}{\partial x}\partial_{x'} + \frac{\partial t'}{\partial x}\partial_{t'}
= \gamma\partial_{x'} - \gamma v \partial_{t'}\\
\partial_t &= \frac{\partial x'}{\partial t}\partial_{x'} + \frac{\partial t'}{\partial t}\partial_{t'}
= -\gamma v\partial_{x'} + \gamma\partial_{t'}.
\end{align}
Now compute the second $x$–derivative:
\begin{align}
\partial_x^2
&= (\gamma\partial_{x'} - \gamma v \partial_{t'})^2 \\
&= \gamma^2\big(\partial_{x'}^2 - 2v\partial_{x'}\partial_{t'} + v^2\partial_{t'}^2\big).
\end{align}
Therefore the Laplacian in the unprimed frame becomes
\begin{align}
\nabla^2
&= \partial_x^2 + \partial_y^2 + \partial_z^2 \\
&= \gamma^2\big(\partial_{x'}^2 - 2v\partial_{x'}\partial_{t'} + v^2\partial_{t'}^2\big)
+ \partial_{y'}^2 + \partial_{z'}^2.
  \end{align}

Compare this to the primed-frame spatial Laplacian,
\begin{equation}
\nabla'^2 = \partial_{x'}^2 + \partial_{y'}^2 + \partial_{z'}^2.
\end{equation}
Clearly $\nabla^2\neq\nabla'^2$ in general: extra mixed and time–derivative terms (those containing $\partial_{t'}$ and $\partial_{x'}\partial_{t'}$) appear when expressing $\nabla^2$ in primed coordinates unless $v=0$ or the field is time-independent in the primed frame. Consequently, applying the transformation to Poisson's equation gives
\begin{equation}
\nabla'^2\Phi + \big(\gamma^2-1\big)\partial_{x'}^2\Phi
-2\gamma^2 v\partial_{x'}\partial_{t'}\Phi + \gamma^2 v^2 \partial_{t'}^2\Phi
= 4\pi G\rho',
\end{equation}
so extra derivative terms of the potential appear on the left-hand side. Since the right-hand side is simply $4\pi G\rho'$ (and $\rho'$ transforms as a scalar density under boosts, not in a way that cancels these additional derivatives), the transformed equation will not in general take the same Poisson form
\begin{equation}
\nabla'^2\Phi(\vec x',t') = 4\pi G\rho(\vec x',t').
\end{equation}

Thus Poisson's equation is \emph{not} invariant under Lorentz transformations: Lorentz covariance instead requires the wave operator (or modifications that include retardation) rather than the purely spatial Laplacian appearing in Poisson's equation.

\section*{Problem 2}
\subsection*{Problem Statement}
There are hypothetical particles calle Tachyons, which always move faster than light.
\begin{itemize}
    \item Show that the four-velocity of Tachyons is space-like.
    \item Also prove that this property is invariant under a Lorentz boos? (Show that if a Tachyon moves faster in one inertial frame, then it moves faster in all inertial frames)
\end{itemize}
\subsection*{Solution}
Assume the Minkowski metric with signature $n_{\mu\nu} = \text{diag}(-1,+1,+1,+1)$. Let $v^\mu$ be a space-like four-velocity (a tachyon $4$-velocity-like vecotr), and let $c^\mu$ be a null (light-like) four-vector.
\begin{quote}
    The usual notion of a normalized four-velocity
    \begin{equation}
        u^\mu u_\mu = -1
    \end{equation}
    does not apply to tachyons.
\end{quote}
By assumption:
\begin{equation}
    v^\mu v_\mu = \eta_{\mu\nu} v^\mu v^\nu > c^\mu c_\mu
\end{equation}
Therefore, since $c^\mu c_\mu = 0$ we've proven that:
\begin{equation}
    v^\mu v_\mu >0
\end{equation}
and thus a space-like four-vector. Now, a Lorentz transformation acts as:
\begin{equation}
    v'^\alpha = \Lambda^\alpha_{\ \ \beta} v^\beta
\end{equation}
with the property:
\begin{equation}
\eta_{\alpha\beta}\Lambda^{\alpha}_{\ \ \mu}\Lambda^{\beta}_{\ \ \nu} = \eta_{\mu\nu}
\end{equation}

By this the Minkwoski norm would be invariant:
\begin{align}
    v'^\alpha v'_\alpha &= \eta_{\alpha\beta}v'^\alpha v'^\beta = \eta_{\alpha\beta}\Lambda^{\alpha}_{\ \ \mu}\Lambda^{\beta}_{\ \ \nu}v^\mu v^\nu \\
    &= \eta_{\mu\nu}v^{\mu}v^\nu = v^\mu v_\mu > 0 \qed
\end{align}
\section*{Problem 3}
\subsection*{Problem Statement}
\begin{itemize}
    \item Show that if two events are time-like separated, there is a Lorentz frame in which they occur at the same point, i.e. at the same spatial coordinate values.
    \item Similarly, show that if two events are space-like separated, there is a Lorentz frame in which they are simultaneous
\end{itemize}
\subsection*{Solution}
Let us imagine two events $a^\mu$ and $b^\mu$. The difference between these two can be shown by:
\begin{equation}
\Delta^\mu = a^\mu - b^\mu
\end{equation}
Now $\Delta^\mu$ is by definition above, a 4-vector. For $|\Delta^\mu|$ there are three scenarios again:

\begin{equation}
\begin{cases}
|\Delta^\mu| > 0, &\text{Space-like}\\
|\Delta^\mu| < 0 , &\text{Time-like}\\
|\Delta^\mu| = 0 , &\text{null}
\end{cases}
\end{equation}
We can in principle choose two coordinates where $\Delta^\mu$ lies on the $x$-axis. Then:

\begin{equation}
\Delta^\mu = (\Delta t , \Delta x , 0,0)
\end{equation}
Assume we're in the $S$ inertial frame, we want to find $S'$ such that:
\begin{equation}
\Delta^{\mu'} = (\Delta t',0,0,0)
\end{equation}
We know the transformation rule between inertial frame are given by the Lorentz transformations therefore:

\begin{equation}
\Delta^{\mu'}= \Lambda^{\mu'}_{\ \ \ \mu} \Delta^{\mu}
\end{equation}
Expanding it means:
\begin{equation}
\begin{cases}
\Delta t' = \gamma(\Delta t  - v \Delta x) \\
\Delta x' = \gamma (\Delta x - v \Delta t)
\end{cases}
\end{equation}
We want $\Delta x' = 0$
Achieving this is possible by finding $v$ of our transformation satisfying:
\begin{align}
\gamma (\Delta x - v\Delta t) &= 0\\
v &= \frac{\Delta x}{\Delta t}
\end{align}

If the events are time-like separated then $\Delta x / \Delta t < 1$. We've shown that there exists a transformation $S \rightarrow S' = S_v$ such that in that frame, the spacial components of the distance between these events would be zero. Interpreting this physically means:

> You moved to the rest frame of the worldline joining the two events.

The crucial condition to be added here is that we need the events be in such distance that $v \not > c$ otherwise the condition only applies when the inertial frame is moving faster than the speed of light.

Following this reasoning we consider that this time $\Delta x / \Delta t >1$. Considering the same equations hold:
\begin{equation}
\begin{cases}
\Delta t' = \gamma(\Delta t  - v \Delta x) \\
\Delta x' = \gamma (\Delta x - v \Delta t)
\end{cases}
\end{equation}
Setting $\Delta t'=0$ would yield:

\begin{align}
\gamma(\Delta t - v\Delta x) &= 0\\
v &= \frac{\Delta t}{\Delta x}
\end{align}

and since $\Delta t / \Delta x < 1$ we're sure that $v < c$, therefore we've shown that there exists a transformation $S \rightarrow S' = S_v$ in which the distance in the final inertial frame have no temporal component.
\section*{Problem 5}
\subsection*{Problem Statement}
Imagine that space (not spacetime) is actually a finite box, or in more sophisticated terms, a three-torus, of size $L$. By this we mean that there is a coordinate system $x^\mu = (t,x,y,z)$ such that every point with coordinates $(t,x,y,z)$ is identified with every point with coordinates $(t,x+L,y,z)$, $(t,x,y+L,z$), and $(t,x,y,z+L)$. Note that the time coordinate is the same. Now consider two observers; observer $A$ is at rest in this coordinate system (constant spatial coordinates), while observer $B$ moves in the $x$-direction with constant velocity $v$. $A$ and $B$ begins at the same events, and while $A$ remains still, $B$ moves once around the universe and comes back to intersect the world-line of $A$ without ever having to accelerate (since the universe is periodic).
\begin{itemize}
\item What are the relative proper times experienced in this interval by $A$ and $B$?
\item Is this consistent with your understanding of Lorentz invariance?
\end{itemize}
\subsection*{Solution}
Space is a $3$-torus of circumference $L$ in each spatial direction, so the identification is:
\begin{equation}
(t,x,y,z)\sim (t,x+L,y,z)
\end{equation}
Time coordinate $t$ is the same at identified points.
Observer $A$ sits at fixed spatial coordinates (so $A$ is at rest in the coordinates that define the identifications). Observer $B$ moves at constant velocity $v$ in the $x$-direction and, because of the periodic identification, returns to intersect $A$'s world-line after one circuit without any acceleration.

We use units with $c=1$ and $0< |v| < 1$.

Along $B$'s path in the given coordinates, the spatial distance travelled in the $x$-direction (from starting event to reunion event) is exactly $\Delta x = L$. Since $B$ moves with velocity $v$ in these coordinates,

\begin{equation}
v = \frac{\Delta x}{\Delta t} \Rightarrow \Delta t = \frac{\Delta x}{v}=\frac{L}{v}
\end{equation}

For observer $A$ (world-line: fixed $x,y,z$):
\begin{equation}
d\tau_A^2 = -ds^2 = dt^2 \Rightarrow \Delta \tau_A = \Delta t = \frac{L}{v}
\end{equation}
For observer $B$ (constant velocity $v$):
The line element for a straight inertial world-line with coordinate velocity $v$ is
\addsvg[0.5]{./spacetime.svg}{spt}{}
\begin{equation}
d\tau_B = dt\sqrt{1-v^2}=\frac{dt}{\gamma}.
\end{equation}
Integrating over the trip,
\begin{equation}
\Delta \tau_B = \Delta t \sqrt{1-v^2}=\frac{\Delta t}{\gamma}=\frac{L}{v\gamma},
\end{equation}
where $\gamma=1/\sqrt{1-v^2}$. Thus
\begin{equation}
\boxed{\Delta \tau_A = \frac{L}{v},  \ \ \Delta\tau_B = \frac{L}{v\gamma}}
\end{equation}
So $B$ ages less: $\Delta\tau_B = \Delta\tau_A/\gamma$. Locally the metric is Minkowski and Lorentz invariance holds: time dilation formula and proper-time calculation used above are standard. The identification $x\sim x+L$ at the same coordinate time $t$ singles out the coordinate frame in which the identifications are purely spatial.
\addsvg[0.5]{./time-dialation.svg}{td}{}
That frame is therefore preferred by the global structure of this spacetime.

If you boost to another inertial frame, the identification no longer connects purely spatial points but connects events separated in time as well. Under a boost with velocity $v$ the identification vector $(\Delta t, \Delta x) = (0,L)$ transforms to:
\begin{equation}
\Delta t' = \gamma(0-vL) = -\gamma vL , \ \ \Delta x' = \gamma L,
\end{equation}
so the identification links events with a nonzero time difference in the primed frame. In other words, the periodic spatial identification breaks global Lorentz symmetry even though the local metric is Lorentzian.

Because of the global asymmetry, the two inertial, non-accelerating world-lines $A$ and $B$ are not equivalent global paths between the same pair of events.
\addsvg[0.5]{./different-homotopy-classes-on-the-torus.svg}{homtop}{}
Topologically, a torus is a closed surface defined as the product of two circles. The moving body ($B$) and the stationary one $A$ are not equal because they lie on different classes of paths (different homotopy classes). Proper time compares lengths of time-like curves
Therefore the result $\Delta\tau_B < \Delta\tau_A$ is fully consistent with special relativity applied locally; the asymmetry arises from the global topology which selects a preferred frame for the identifications.

\section*{Probllem 6}
\subsection*{Problem Statement}
Three events $A$,$B$,$C$ are seen by observer $\mathcal{O}$ to occur in the order $ABC$. Another observer $\bar{\mathcal{O}}$, sees the events to occur in the order $CBA$. Is it possible that a third observer sees the events in the order $ACB$? Support you conclusion by drawing a spacetime diagram.
\subsection*{Solution}
And
\section*{Problem 7}
\subsection*{Problem Statment}
Particle physicists are so used to setting $c=1$ that they measure mass in units of energy. In particular, they tend to use electron-volts ($1\text{eV} = 1.6 \times10^{-12}\text{erg} = 1.8\times10^{-33}\text{gr}$), or, more commonly $\text{keV}$, $\text{MeV}$ and $\text{GeV}$ ($10^3 \text{eV}, 10^6\text{eV}, 10^9\text{eV}$, respectively).

The muon has been measured to have a mass of $0.106\text{GeV}$ and a rest frame life-time of $2.19\times 10^{-6}$ seconds. Imagine that such a muon is moving in the circular storage ring of a particle accelerator, $1$ kilometer in diameter, such that the muon's total energy is $1 \text{TeV}$. How long would it appear to live from the experimenter's point of view? How many radians would it travel around the ring?

\subsection*{Solution}
We work in units with $c=1$. The relativistic gamma factor is:
\begin{equation}
    \gamma=\frac{E}{m} = \frac{1000\text{GeV}}{0.106\text{GeV}} \approx 9.4\times 10^3
\end{equation}
One can calculate the dilated lifetime seen by the experimenter as below:
\begin{equation}
    \tau_{\text{lab}} = \gamma \tau_0 = (9.4 \times 10^3)\times(2.19\times 10^{-6}\text{s})\approx 2\time 10^{-2}
\end{equation}
The muon speed folows from:
\begin{equation}
    \beta = \sqrt{1-\frac{1}{\gamma^2}}
\end{equation}
The storage ring has diameter $1$ kilometer, so radius $r=500$ meter and circumference:
\begin{equation}
    C = 2\pi r = 3141 m
\end{equation}
cal calculating the angle traveled:
\begin{equation}
    \theta = \frac{d}{r} = \frac{v \tau_{\text{lab}}} = 1.23\times 10^4\text{rad}
\end{equation}
\section*{Problem 8}
\subsection*{Problem Statement}
In Euclidean three-space, let $p$ be the point with coordinates $(x,y,z) = (1,0,-1)$. Consider the following curves that pass through $p$:

\begin{align}
x^i(\lambda) &= (\lambda, (\lambda - 1)^2, -\lambda) \\
x^i (\mu) &= (\cos\mu,\sin\mu,\mu-1)\\
x^i(\sigma)&=(\sigma^2,\sigma^3+\sigma^2,\sigma)\\
\end{align}
\begin{itemize}
\item Calculate the components of the tangent vectors to these curves at $p$ in the coordinate basis $\{\partial_x,\partial_y,\partial_z\}$
\item Let $f=x^2 + y^2 - yz$. Calculate $\frac{df}{d\lambda},\frac{df}{d\mu},\frac{df}{d\sigma}$
\end{itemize}
\subsection*{Solution}

The tangent vector to a curve at parameter value $t$ is given by $\frac{dx^i}{dt}$ evaluated at that parameter.

\subsubsection*{Tangent vector for curve 1 (at $\lambda = 1$):}

\begin{align*}
\frac{dx^i}{d\lambda} &= \left(\frac{d\lambda}{d\lambda}, \frac{d(\lambda-1)^2}{d\lambda}, \frac{d(-\lambda)}{d\lambda}\right) \\
&= (1, 2(\lambda-1), -1)
\end{align*}

At $\lambda = 1$:
\begin{equation*}
\boxed{\left.\frac{dx^i}{d\lambda}\right|_{\lambda=1} = (1, 0, -1)}
\end{equation*}

\subsubsection*{Tangent vector for curve 2 (at $\mu = 0$):}

\begin{align*}
\frac{dx^i}{d\mu} &= \left(\frac{d(\cos\mu)}{d\mu}, \frac{d(\sin\mu)}{d\mu}, \frac{d(\mu-1)}{d\mu}\right) \\
&= (-\sin\mu, \cos\mu, 1)
\end{align*}

At $\mu = 0$:
\begin{equation*}
\boxed{\left.\frac{dx^i}{d\mu}\right|_{\mu=0} = (0, 1, 1)}
\end{equation*}

\subsubsection*{Tangent vector for curve 3 (at $\sigma = -1$):}

\begin{align*}
\frac{dx^i}{d\sigma} &= \left(\frac{d(\sigma^2)}{d\sigma}, \frac{d(\sigma^3+\sigma^2)}{d\sigma}, \frac{d\sigma}{d\sigma}\right) \\
&= (2\sigma, 3\sigma^2 + 2\sigma, 1)
\end{align*}

At $\sigma = -1$:
\begin{equation*}
\boxed{\left.\frac{dx^i}{d\sigma}\right|_{\sigma=-1} = (-2, 1, 1)}
\end{equation*}

\subsection*{Part 2: Derivatives of $f$ along the curves}

Given $f = x^2 + y^2 - yz$, we use the chain rule:
\begin{equation*}
\frac{df}{dt} = \frac{\partial f}{\partial x}\frac{dx}{dt} + \frac{\partial f}{\partial y}\frac{dy}{dt} + \frac{\partial f}{\partial z}\frac{dz}{dt}
\end{equation*}

First, compute the gradient of $f$:
\begin{equation*}
\nabla f = (2x, 2y - z, -y)
\end{equation*}

At $p = (1, 0, -1)$:
\begin{equation*}
\nabla f|_p = (2, 1, 0)
\end{equation*}

\subsubsection*{For curve 1:}
\begin{align*}
\frac{df}{d\lambda} &= \nabla f \cdot \frac{dx^i}{d\lambda} \\
&= (2, 1, 0) \cdot (1, 0, -1) \\
&= 2 + 0 + 0
\end{align*}
\begin{equation*}
\boxed{\frac{df}{d\lambda} = 2}
\end{equation*}

\subsubsection*{For curve 2:}
\begin{align*}
\frac{df}{d\mu} &= \nabla f \cdot \frac{dx^i}{d\mu} \\
&= (2, 1, 0) \cdot (0, 1, 1) \\
&= 0 + 1 + 0
\end{align*}
\begin{equation*}
\boxed{\frac{df}{d\mu} = 1}
\end{equation*}

\subsubsection*{For curve 3:}
\begin{align*}
\frac{df}{d\sigma} &= \nabla f \cdot \frac{dx^i}{d\sigma} \\
&= (2, 1, 0) \cdot (-2, 1, 1) \\
&= -4 + 1 + 0
\end{align*}
\begin{equation*}
\boxed{\frac{df}{d\sigma} = -3}
\end{equation*}
\end{document}